% !TEX root = ../main.tex
\chapter{Introduction and Research Problem}
\label{ch:introduction}

Decentralized Finance (DeFi) has transformed how capital is allocated on public blockchains, enabling lending, borrowing, and leveraged trading without traditional financial intermediaries (Sch\"{a}r, 2021; Harvey et al., 2021). Yet most credit-like interactions remain highly collateralized because protocols cannot rely on off-chain identity systems or conventional credit histories (Packin \& Lev-Aretz, 2024; Gudgeon et al., 2020). In a pseudonymous environment, the operational challenge is to allocate capital safely: distinguishing wallets that manage leveraged exposure reliably from those that exhibit repeated losses, liquidations, or adversarial behavior.

This dissertation addresses a concrete research problem within that setting: whether an on-chain reputation score built from ERC-20 \texttt{approve}/\texttt{permit} allowances---implemented as \textbf{EndorseRank}---can reduce computational cost, represent endorsement semantics more faithfully, and yield a distinct yet credible alignment profile compared with transaction-based Adaptive Weighted PageRank (AWP) on \textbf{Arbitrum One}. The primary empirical comparison is \textbf{EndorseRank versus AWP} on an identical wallet sample, validated with five proxy families and GMX V2 perpetual swap outcomes following Do et al.'s (2023) rank-correlation framework.

In this study, \textbf{on-chain reputation score} computation refers to deriving a quantitative wallet score from observable on-chain behavior and relationships without off-chain identity (Do et al., 2023). EndorseRank produces such a score from allowance-based endorsement edges. The remainder of this chapter situates the problem historically and economically, motivates allowance-based reputation, states the research problem and questions, defines key constructs, and outlines contributions, scope, and the dissertation roadmap.

\dissertationSubheading[sec:background]{Background: DeFi, collateralization, and pseudonymity}

Public blockchains introduced programmable settlement and open participation without a central operator (Nakamoto, 2008; Wood, 2014). Smart contracts encode agreement terms as self-executing code, enabling financial applications that settle on-chain rather than through custodial intermediaries (Buterin, 2014; Szabo, 1997). DeFi protocols compose these contracts into markets for lending, automated market making, and perpetual futures (Sch\"{a}r, 2021; Werner et al., 2022). On Arbitrum One and similar Layer-2 networks, activity density is high enough that wallet-level behavioral signals---transfers, approvals, and protocol events---can be observed at scale.

Despite this transparency, DeFi credit markets remain structurally over-collateralized. Protocols typically require borrowers or traders to post collateral exceeding the value of their exposure because counterparties are pseudonymous addresses rather than identified legal persons (Gudgeon et al., 2020; Harvey et al., 2021). Pseudonymity is a design feature of public ledgers: anyone may create addresses at negligible cost, and a single economic actor may control many wallets (Douceur, 2002). Traditional credit scoring depends on identity-linked histories that do not transfer cleanly to this setting (Packin \& Lev-Aretz, 2024). Consequently, risk management defaults to collateral buffers, liquidation engines, and protocol-level parameters rather than individualized reputation.

Collateralization protects protocols but imposes capital inefficiency. Traders and borrowers must lock assets that could otherwise be deployed elsewhere; protocols forgo under-collateralized products that traditional finance offers to borrowers with established credit. Researchers have therefore explored on-chain signals as substitutes for off-chain identity: transfer-graph centrality (Do et al., 2023), network risk propagation (Lin et al., 2024), prospect-theory and blockchain-based credit models (Hassija et al., 2020), and machine-learning pipelines on lending protocols (Kotb, 2023). These approaches share a premise: public ledger data encode economically meaningful relationships that can inform risk assessment without deanonymizing users.

Graph ranking methods are a natural fit for this premise. PageRank and its variants propagate influence through directed edges, ranking nodes by recursive endorsement from other nodes (Page et al., 1998; Brin \& Page, 1998; Langville \& Meyer, 2006). Early web-trust systems such as EigenTrust and TrustRank applied similar ideas to peer-to-peer and web spam settings (Kamvar et al., 2003; Gy\"{o}ngyi et al., 2004). Do et al.\ (2023) adapt PageRank to Ethereum transfer graphs with value-weighted, time-decayed edges---Adaptive Weighted PageRank (AWP)---and validate scores against inbound transfer volume as a domain-intuition proxy for economic importance. That line of work establishes transfer-graph reputation as a credible baseline; it does not, however, model explicit spending authorization, which is a distinct on-chain act recorded by ERC-20 allowances (Ethereum Improvement Proposals, 2015, 2020).

\dissertationSubheading[sec:motivation]{Motivation for on-chain reputation}

Three practical pressures motivate on-chain reputation research in leveraged DeFi. First, \textbf{capital efficiency}: if protocols can identify wallets with historically reliable risk management, they may safely reduce collateral requirements or prioritize liquidations, improving utilization without abandoning on-chain settlement (Sch\"{a}r, 2021; Gudgeon et al., 2020). Second, \textbf{auditability}: unlike opaque off-chain credit scores, graph-based scores derived from public events can be recomputed and inspected, addressing concerns about black-box decentralized credit scoring (Packin \& Lev-Aretz, 2024). Third, \textbf{operational cost}: reputation systems that must refresh frequently---for example, before each leverage adjustment---cannot afford expensive historical transfer aggregation on every update.

Transfer-based AWP addresses the first two pressures but is costly on the third. Constructing $G_{\text{transfer}}$ requires aggregating transfer history, applying logistic time-decay, and running PageRank on a dense edge set (Do et al., 2023). On the matched cohort studied here ($n=5{,}521$ wallets), the AWP subgraph contains $137{,}087$ edges and requires $18.711$\,s of PageRank wall-clock time. Allowance graphs offer a structurally different alternative. An ERC-20 \texttt{approve} or EIP-2612 \texttt{permit} records how much a token owner authorizes a spender to transfer (Ethereum Improvement Proposals, 2015, 2020). The latest allowance per owner--spender pair is a snapshot of current authorization intent: newer approvals override older ones, so temporal decay is unnecessary. The resulting endorsement graph $G_{\text{approve}}$ is sparse---$14{,}727$ edges on the same cohort---and PageRank completes in $2.105$\,s, approximately $8.9\times$ faster than AWP.

Beyond efficiency, allowances carry distinct semantics. A transfer may reflect routing, market making, or passive receipt; an allowance is an explicit delegation of spending authority. Treating allowances as endorsement edges aligns PageRank's citation metaphor (Page et al., 1998) with a trust-like act that protocols already rely on for token interactions. Whether that semantic difference produces rankings that are merely different, or differently aligned with economically relevant outcomes, is an empirical question. This dissertation answers that question on Arbitrum One using GMX V2 perpetual trading outcomes as a dense validation frame (GMX, n.d.), because under-collateralized lending default labels remain sparse on this chain.

\dissertationSubheading[sec:research-problem]{Research problem}

Existing Adaptive Weighted PageRank (AWP) methods infer trust from transfer volume and logistic time-decay over transaction history (Do et al., 2023). While economically interpretable, they (i) require expensive historical transfer aggregation, (ii) treat inbound transfer volume as a proxy for a ``valuable wallet,'' and (iii) do not model explicit spending authorization. EndorseRank instead uses the latest allowance per owner--spender pair as edge strength on a directed endorsement graph, yielding a sparser graph and a semantically distinct trust signal without temporal decay weighting---integrating ERC-20 allowance edges into PageRank for enhanced on-chain wallet reputation scoring.

The research problem is therefore comparative and evaluative: on an identical Arbitrum One wallet sample with GMX V2 perpetual activity, does EndorseRank produce wallet orderings that are construct-distinct from AWP, align credibly with domain-intuition proxies (including allowance-specific and trading-success families), and do so at substantially lower computational cost? The primary comparison is EndorseRank versus AWP; a supplementary outcome-native reference (GF-PR) is introduced later to bound interpretation of trading-outcome proxies, not to replace either social method.

Table~\ref{tab:primary-methods-preview} summarizes the primary comparison. Table~\ref{tab:hrq-claim-map} summarizes hypotheses (H1--H3), research questions (RQ1--RQ4), supplementary question SRQ1, and Chapter~4 empirical claims (C1--C4, SC1).

\begin{table}[htbp]
\centering
\caption{Primary comparison: EndorseRank versus AWP (matched wallet cohort, $n=5{,}521$).}
\label{tab:primary-methods-preview}
\begin{tabular}{p{2.4cm}p{3.4cm}p{7.4cm}}
\toprule
Method & Input graph & Role \\
\midrule
EndorseRank & Latest ERC-20 allowances ($G_{\text{approve}}$) & Primary construct: allowance delegation; target of this dissertation \\
AWP & Time-decayed transfers ($G_{\text{transfer}}$) & Primary baseline: Do et al.\ transfer-graph reproduction \\
\bottomrule
\end{tabular}
\end{table}

To avoid ambiguity, the key concepts used throughout this dissertation are defined as follows:

\begin{itemize}
\item \textbf{DeFi (Decentralized Finance):} Financial services---including lending, borrowing, and perpetual trading---implemented via smart contracts on public blockchains, typically without a central operator controlling custody or settlement (Sch\"{a}r, 2021; Harvey et al., 2021). DeFi protocols compose open interfaces so that positions, liquidations, and token movements are observable on-chain.
\item \textbf{Smart contract:} Self-executing program code deployed on a blockchain that enforces agreement terms without a trusted intermediary (Buterin, 2014; Szabo, 1997). In this dissertation, smart contracts are the source of ERC-20 allowance logs, transfer logs, and GMX V2 position events.
\item \textbf{Wallet / Address:} A blockchain account that holds assets and interacts with smart contracts; the unit of analysis for reputation scoring (Buterin, 2014; Ethereum Improvement Proposals, 2015). A wallet is identified by its address; no off-chain identity is assumed or recovered.
\item \textbf{On-chain reputation score:} A numerical indicator derived from on-chain signals that represents a wallet's relative standing for financial interactions under pseudonymity (Do et al., 2023). Scores are ordinal for ranking purposes; absolute magnitudes are not interpreted as credit limits.
\item \textbf{On-chain credit risk:} Observable adverse outcomes such as position liquidation or repeated trading losses on public blockchains (Lin et al., 2024). This dissertation uses GMX V2 perpetual outcomes as domain-intuition proxies for risk-relevant behavior, not as loan-repayment ground truth.
\item \textbf{GMX V2 perpetual swap position:} A collateralized long or short exposure on GMX V2, opened via \textbf{PositionIncrease} and closed via \textbf{PositionDecrease} events on Arbitrum One (GMX, n.d.; Sch\"{a}r, 2021). Realized profit-and-loss from closes supplies trading-success and inverse-risk proxies.
\item \textbf{PageRank:} A graph ranking method that propagates influence through directed edges under a damping factor (Page et al., 1998; Brin \& Page, 1998; Langville \& Meyer, 2006). \textbf{EndorseRank} applies PageRank to a latest-allowance endorsement graph: edge strength equals the current allowance per owner--spender pair, with no temporal decay. \textbf{Adaptive Weighted PageRank (AWP)} applies PageRank to transfer graphs with edge weights given by transfer value and logistic time-decay (Do et al., 2023).
\item \textbf{ERC-20 allowance / Approve / Permit:} Mechanisms recording how much a token owner permits a spender to transfer on the owner's behalf (Ethereum Improvement Proposals, 2015, 2020). \texttt{Approve} is a standard on-chain call; EIP-2612 \texttt{permit} authorizes the same state change via signature. EndorseRank uses the latest non-zero allowance snapshot per owner--spender pair.
\item \textbf{Endorsement graph:} A directed graph in which edges represent explicit delegation of spending authority rather than mere asset transfers (Do et al., 2023). In EndorseRank, an edge $u \to v$ means wallet $u$ currently authorizes wallet or contract $v$ to spend tokens owned by $u$.
\item \textbf{Domain alignment:} The degree to which a reputation ranking agrees with an observable proxy ranking, assessed by Spearman's $\rho$ and Kendall's $\tau$ (Spearman, 1904; Kendall, 1938; Cronbach \& Meehl, 1955; Do et al., 2023). Alignment is construct-specific: a method may align strongly with its primary graph construct and weakly with an orthogonal outcome family.
\end{itemize}

\dissertationSubheading[sec:supplementary-extension]{Supplementary extension (GF-PR)}

During empirical analysis, \textbf{GF-PR} (GainFlow PageRank) was added as a supplementary outcome-native reference (Section~\ref{sub:gf-pr}). GF-PR builds a directed star graph from the same GMX profit-and-loss stream used to define validation proxies. It ranks wallets on synthetic POOL/SINK flow rather than social wallet-to-wallet edges, setting an upper bound on trading-outcome proxy alignment for diagnostic interpretation---not a deployable social reputation score.

\begin{table}[htbp]
\centering
\caption{Supplementary extension added during empirical work (not part of the primary EndorseRank--AWP research question).}
\label{tab:supplementary-gf-pr}
\begin{tabular}{p{2.4cm}p{3.4cm}p{7.4cm}}
\toprule
Method & Input graph & Role \\
\midrule
GF-PR & GMX PnL star (POOL$\to$wallet, wallet$\to$SINK) & Supplementary ceiling reference; construct overlap on trading-outcome proxies expected \\
\bottomrule
\end{tabular}
\end{table}

\begin{table}[htbp]
\centering
\caption{Summary of hypotheses, research questions, and empirical claims.}
\label{tab:hrq-claim-map}
\begin{tabular}{p{1.5cm}p{6.0cm}p{5.6cm}}
\toprule
Item & Addresses & Reported in \\
\midrule
H1 / RQ1 & Do EndorseRank and AWP yield distinct wallet orderings? & \S\ref{sec:rank-divergence}; inter-method $\tau=0.316$ \\
H2 / RQ2 & Which social method aligns more strongly with which proxy families? & Tables~\ref{tab:alignment-three-methods}--\ref{tab:alignment-gmx}; Claims C2--C3 \\
H3 / RQ3 & Is EndorseRank faster than AWP at equal $n$? & Tables~\ref{tab:benchmark-runtime}--\ref{tab:benchmark-scaling}; Claims C1, C4 \\
RQ4 & Transfer vs.\ trading-success validation frame for social methods & Table~\ref{tab:alignment-three-methods} \\
SRQ1 & What does GF-PR ceiling imply for trading-outcome proxy interpretation? & \S\ref{sec:gf-pr-ceiling}; Claim SC1 \\
C1 & EndorseRank vs.\ AWP runtime and graph size & Table~\ref{tab:benchmark-runtime} \\
C2 & Credible EndorseRank--AWP alignment on selected proxies & Tables~\ref{tab:alignment-transfer}--\ref{tab:alignment-gmx} \\
C3 & Construct-specific winners (EndorseRank vs.\ AWP) & Table~\ref{tab:alignment-three-methods} \\
C4 & Efficiency--alignment trade-off for EndorseRank & Tables~\ref{tab:benchmark-runtime}--\ref{tab:benchmark-scaling} \\
SC1 & Supplementary GF-PR construct-overlap diagnostic & \S\ref{sec:gf-pr-ceiling} \\
\bottomrule
\end{tabular}
\end{table}

\dissertationSubheading[sec:research-questions]{Research questions}

This study tests three primary hypotheses: \textbf{H1}---EndorseRank and AWP yield distinct wallet orderings on the same cohort; \textbf{H2}---the two social methods differ in construct-specific proxy alignment; and \textbf{H3}---EndorseRank is computationally more efficient than AWP at equal $n$. The following research questions guide the empirical evaluation summarized in Table~\ref{tab:hrq-claim-map}. RQ1--RQ4 form the primary EndorseRank--AWP comparison; SRQ1 addresses the supplementary GF-PR extension.

\begin{enumerate}
\item[\textbf{RQ1.}] Do EndorseRank and AWP produce meaningfully different wallet orderings on an identical Arbitrum wallet sample?
\item[\textbf{RQ2.}] Between EndorseRank and AWP, which social method aligns more strongly with each of the five proxy families (Spearman's $\rho$, Kendall's $\tau$)?
\item[\textbf{RQ3.}] Is EndorseRank computationally more efficient than AWP on the same matched sample (runtime, memory, iterations)?
\item[\textbf{RQ4.}] For social reputation methods, which validation frame yields stronger rank alignment on this chain: Do et al.\ transfer proxies (in-degree/in-value) or trading-success proxies (from GMX V2 perpetual closes)?
\item[\textbf{SRQ1.}] What does the supplementary GF-PR outcome-native ceiling imply for interpreting trading-outcome proxy alignment of EndorseRank and AWP?
\end{enumerate}

RQ1 tests whether allowance and transfer graphs induce distinct centrality structures, operationalized by inter-method Kendall $\tau$ between score vectors. RQ2 evaluates construct-specific alignment across transfer, allowance, inverse-risk, Sybil-adjusted stability, and trading-success proxy families. RQ3 compares PageRank wall-clock time, memory, and iterations on identical wallet sets, including scaling on an expanded wallet pool. RQ4 asks which validation frame---transfer centrality or trading success---is more informative for social reputation methods on this sample. SRQ1 uses GF-PR only to bound how high trading-outcome alignment can rise when ranking directly on outcome flow, clarifying construct overlap (Cronbach \& Meehl, 1955).

\dissertationSubheading[sec:contributions]{Contributions}

This dissertation makes four contributions to on-chain reputation research and DeFi risk assessment.

First, it formulates \textbf{EndorseRank}, an allowance-based endorsement graph model that applies PageRank to latest ERC-20 \texttt{approve}/\texttt{permit} edges. The model treats current spending authorization as endorsement strength, eliminates temporal decay by relying on allowance override semantics, and produces a sparse graph suitable for frequent refresh.

Second, it conducts a \textbf{five-proxy domain alignment evaluation} comparing EndorseRank and AWP on leveraged DeFi activity. Proxies span transfer centrality (Do et al., 2023), allowance delegation, GMX V2 inverse risk and trading success (GMX, n.d.), and Sybil-adjusted stability. Alignment is reported with Spearman's $\rho$ and Kendall's $\tau$ (Spearman, 1904; Kendall, 1938), enabling construct-specific interpretation rather than a single global winner.

Third, it releases an \textbf{open reproducible evaluation pipeline} that links BigQuery extraction, graph construction, PageRank scoring, proxy computation, and LaTeX result tables to a machine-readable summary (Appendix~\ref{ch:appendix-eval}). The pipeline supports independent verification of claims C1--C4 and SC1.

Fourth, as a supplementary contribution added during empirical work, it provides a \textbf{GF-PR ceiling analysis} that clarifies construct overlap on trading-outcome proxies for the matched wallet cohort ($n=5{,}521$). GF-PR is not proposed as a deployable social reputation method; it bounds interpretation of high $\tau$ on outcome-native families (Claim~SC1).

Empirically, the primary comparison yields inter-method $\tau=0.316$, EndorseRank allowance $\tau=0.355$, AWP transfer $\tau=0.534$, and an efficiency advantage of $2.105$\,s versus $18.711$\,s ($\approx 8.9\times$) on the matched cohort, with pool scaling of $\approx 10.4\times$ at $n=10{,}000$ and $\approx 9.4\times$ at $N=31{,}612$. These results support an efficiency--alignment trade-off rather than universal dominance of either social method.

\dissertationSubheading[sec:scope-limitations]{Scope and limitations}

This dissertation focuses on Arbitrum One and GMX V2 perpetual trading over the observation window 2025-12-01 to 2026-05-31. The primary empirical sample is a matched wallet cohort of $n=5{,}521$ addresses with at least three GMX \texttt{PositionDecrease} events and non-zero reputation subgraph connectivity, comprising $365{,}488$ decoded closes, $14{,}727$ EndorseRank edges, and $137{,}087$ AWP edges. Runtime scaling is reported on an expanded wallet pool ($N=31{,}612$). Robustness checks examine damping sensitivity, top-token subgraphs, and sample-size sensitivity.

Unsecured lending default labels are out of scope: under-collateralized credit events are too sparse on this chain for reliable ground-truth validation. Constructs and ethics are defined in Chapters~\ref{ch:introduction} and~\ref{ch:methodology}. Among social methods, transfer proxies align more strongly than trading-success proxies on this sample (RQ4); this does not diminish EndorseRank's allowance-specific validity or runtime advantage. GF-PR results are reported as supplementary diagnostics (SRQ1), not as a third peer reputation method. Extended baselines archived for future comparison are excluded from the main narrative (Appendix~\ref{ch:appendix-eval}).

Generalization beyond Arbitrum One, GMX V2, and the six-month window is not claimed. Token-decimal normalization, price oracles, and permit-versus-approve coverage may affect edge weights. Adversarial manipulation of allowances or transfers remains a threat discussed in Chapter~\ref{ch:discussion}; Sybil-adjusted proxies partially address but do not eliminate that risk (Douceur, 2002).

\dissertationSubheading[sec:roadmap]{Dissertation roadmap}

The dissertation proceeds as follows.

Chapter~\ref{ch:literature} reviews graph-based on-chain reputation, ERC-20 allowance semantics, and GMX V2 trading-success validation. It develops the theoretical framework---hyperlink propagation, domain alignment, and computational complexity---that motivates the primary EndorseRank versus AWP comparison, and situates GF-PR as a supplementary outcome-native ceiling.

Chapter~\ref{ch:methodology} specifies data sources, sampling, graph construction for EndorseRank and AWP, the five-proxy evaluation design, computational benchmarks, robustness checks, reproducibility practices, and ethics. It defines the matched wallet cohort, the expanded wallet pool used for scaling, and the operational procedures that produce the Chapter~\ref{ch:results} tables.

Chapter~\ref{ch:results} reports empirical results: cohort characteristics, runtime and scaling benchmarks (Claim~C1), robustness checks, social-method alignment across five proxy families (Claims~C2--C3), rank divergence (H1/RQ1), and the supplementary GF-PR ceiling analysis (SRQ1/SC1). Primary claims C1--C4 and supplementary claim SC1 are stated explicitly.

Chapter~\ref{ch:discussion} interprets findings by research question, develops theoretical and practical implications, formulates a Sybil attack cost model on allowance graphs, discusses regulatory and ethical considerations, and evaluates threats to validity and limitations.

Chapter~\ref{ch:conclusion} synthesizes the contribution to knowledge, summarizes findings against RQ1--RQ4 and SRQ1, and outlines future work on cross-protocol validation, adversarial robustness, hybrid methods, and live oracle integration.
